% ============================================================
%  SECCIÓN 3: ENTENDIMIENTO DE LA PLANTA Y MODELADO MATEMÁTICO
% ============================================================
\section{Entendimiento de la planta y Modelado Matem\'atico}

\subsection{Ecuaci\'on Diferencial No Lineal}
A partir del diagrama de cuerpo libre del aerop\'endulo en rotaci\'on plana alrededor del eje fijo del rodamiento, la sumatoria de momentos satisface:
\begin{equation}
  \sum \tau_O = I \ddot{\theta}(t)
\end{equation}

Considerando el torque generado por el empuje aerodin\'amico del motor $F(t) = K_{uf} u(t)$ a distancia radial $r$, el momento restaurador gravitacional con centro de masa a distancia $c$ y el amortiguamiento viscoso equivalente $B r^2$:
\begin{equation}
  I \ddot{\theta}(t) = r F(t) - B r^2 \dot{\theta}(t) - m g c \cos(\theta(t))
  \label{eq:dinamica_nolineal}
\end{equation}
donde $\theta(t)$ representa la posici\'on angular, $I$ el momento de inercia respecto al pivote, $m$ la masa oscilante, $g$ la aceleraci\'on de la gravedad y $u(t)$ la se\~nal de mando.

\subsection{Puntos de Equilibrio Estacionario}
En r\'egimen permanente estacionario ($\ddot{\theta}^* = 0$, $\dot{\theta}^* = 0$), la fuerza requerida para sostener un \'angulo de operaci\'on $\theta^*$ satisface:
\begin{equation}
  F^* = \frac{m g c}{r} \cos(\theta^*)
\end{equation}

\subsection{Linealizaci\'on Alrededor del Punto de Operaci\'on}
Dado el comportamiento fuertemente no lineal del t\'ermino trigonom\'etrico, se plantea una linealizaci\'on mediante aproximaci\'on por serie de Taylor alrededor de la condici\'on de operaci\'on $(\theta^*, \dot{\theta}^* = 0, F^*)$.

Definiendo las variables de desviaci\'on incremental:
\begin{equation}
  \Delta \theta(t) = \theta(t) - \theta^*, \qquad \Delta F(t) = F(t) - F^*
\end{equation}
la din\'amica linealizada resultante adquiere la forma:
\begin{equation}
  I \Delta \ddot{\theta}(t) + B r^2 \Delta \dot{\theta}(t) - m g c \sin(\theta^*) \Delta \theta(t) = r \Delta F(t)
  \label{eq:dinamica_linealizada}
\end{equation}

\subsection{Representaci\'on del Modelo (a definir con el equipo)}
A partir de la ecuaci\'on linealizada (\ref{eq:dinamica_linealizada}), se estructuran las v\'{\i}as de representaci\'on para la s\'{\i}ntesis del controlador:
\begin{itemize}[leftmargin=*]
  \item \textbf{Funci\'on de transferencia incremental $G(s)$:}
  \begin{equation}
    G(s) = \frac{\Delta \Theta(s)}{\Delta F(s)} = \frac{\dfrac{r}{I}}{s^2 + \dfrac{B r^2}{I} s - \dfrac{m g c \sin(\theta^*)}{I}}
  \end{equation}
  \item \textbf{Representaci\'on en espacio de estados:}
  Formulaci\'on matricial $\Delta \dot{\mathbf{x}} = \mathbf{A} \Delta \mathbf{x} + \mathbf{B} \Delta u$ para t\'ecnicas de control multivariable o moderno.
\end{itemize}
